%% ma: L12-B14-0017
%% lop: 12
%% chuong: 5
%% bai: 14
%% dang: 12.14.3
%% loai: TN4
%% muc: NB
%% dap_an: "C"
%% nguon: "CĐ2.5.pdf giáo viên gửi 27/10/2026 (M07, Lê Minh Khôi – Nguyễn Anh Khôi, Chương II Chủ đề 5 Tích có hướng của hai vectơ), Câu 1"
%% kien_thuc: "Bài 14 (tích có hướng của hai vectơ)"
%% trang_thai: chinh-thuc
\begin{ex}
Trong không gian $Oxyz$, cho hai vectơ $\vec a=(a_1;a_2;a_3)$ và $\vec b=(b_1;b_2;b_3)$. Tích có hướng của hai vectơ $\vec a$ và $\vec b$ được định nghĩa là
\choice{$\left[\vec a,\vec b\right]=(a_1b_1-a_2b_2;a_2b_2-a_3b_3;a_3b_3-a_1b_1)$}{$\left[\vec a,\vec b\right]=(a_2b_3-a_3b_2;a_1b_3-a_3b_1;a_1b_2-a_2b_1)$}{\True $\left[\vec a,\vec b\right]=(a_2b_3-a_3b_2;a_3b_1-a_1b_3;a_1b_2-a_2b_1)$}{$\left[\vec a,\vec b\right]=(a_1b_3-a_3b_1;a_2b_1-a_1b_2;a_2b_3-a_3b_2)$}
\loigiai{$\left[\vec a,\vec b\right]=\left(\begin{vmatrix}a_2&a_3\\b_2&b_3\end{vmatrix};\begin{vmatrix}a_3&a_1\\b_3&b_1\end{vmatrix};\begin{vmatrix}a_1&a_2\\b_1&b_2\end{vmatrix}\right)=(a_2b_3-a_3b_2;a_3b_1-a_1b_3;a_1b_2-a_2b_1)$.}
\end{ex}
